% Definições básicas do resumo
\begin{adjustwidth}{-1cm}{-1cm}\vspace{\onelineskip}

%% Resumo

\begin{resumo}

O desenvolvimento de plataformas robóticas móveis em ambientes educacionais e de pesquisa demanda arquiteturas flexíveis e acessíveis. Este trabalho propõe a adaptação de um robô aspirador comercial (Roomba® 614 iRobot) em uma plataforma robótica seguidora de linha autônoma. O sistema utiliza um Raspberry Pi 4 Model B, operando com o sistema Ubuntu 22.04, como unidade central de processamento. Essa escolha elimina a necessidade de microcontroladores intermediários e permite a leitura direta de um módulo seguidor de linha de 5 canais (BFD-1000) e a execução do algoritmo de controle Proporcional-Integral-Derivativo (PID) via Edge Computing. Todo o software de controle e comunicação foi desenvolvido de forma modular sobre o ecossistema \textit{Robot Operating System 2} (ROS 2 Humble), adotando a linguagem C++ e Python. A integração destas tecnologias resultou em uma arquitetura de hardware simplificada e robusta, empregando um cabo conversor USB-Serial TTL para mitigar ruídos de comunicação. Os resultados demonstram a validação validação experimental da plataforma em ambiente físico e o potencial da plataforma como um recurso didático avançado para o ensino de cinemática, automação e sistemas embarcados.

\vspace{\onelineskip}
\noindent

\textbf{Palavras-chave}: robótica móvel; seguidor de linha; ROS 2 Humble; Raspberry Pi 4; controle PID.

\end{resumo}
\end{adjustwidth}